Breast ultrasound is a useful adjunct to mammography, but its interpretation is operator dependent and image appearance is affected by speckle, acquisition settings, and lesion context. This dissertation develops a reproducible research prototype for binary classification of breast-ultrasound images using convolutional neural networks. The software integrates contrast-limited adaptive histogram equalisation (CLAHE), aspect-ratio-preserving square padding, transfer-learning backbones (EfficientNet-B0 and ResNet-50), and a custom convolutional baseline. It also provides a FastAPI research demonstration for controlled image inference and visual inspection.

The principal contribution is methodological rather than clinical: a single preprocessing implementation is shared by training, evaluation, command-line inference, and the web service; paired lesion masks are recognised under the naming conventions used by the supplied datasets; and evaluation emits machine-readable metrics, figures, and an audit trail. A read-only data audit records split sizes and flags cross-split subject identifiers before performance is interpreted.

The local OASBUD and BrEaST copies were rebuilt into subject-level partitions and pass the split audit. BUSI remains excluded because its renamed image copy does not retain enough provenance to verify patient separation. A fresh EfficientNet-B0 checkpoint evaluated on 62 held-out BrEaST and OASBUD images produced 67.74\% accuracy, macro F1 0.6774, and ROC-AUC 0.7419. These values describe a local research cohort, not clinical or external validation. Historical headline figures, pseudo-labelling ablations, and noise curves contained in legacy project assets are therefore not treated as empirical findings.

The corrected workflow rebuilds the cohort at subject level, preserves source identifiers and preprocessing manifests, retrains with fixed seeds, and generates tables and figures from held-out predictions. The system remains an educational research prototype, not a diagnostic or clinical decision-support device.
